\section{Introduction}
\label{sec:introduction}

Federated learning (FL) enables many clients to train a shared model without centralizing their private data \citep{mcmahan2017fedavg,bonawitz2017secureagg,kairouz2021federated}. This distributed training paradigm is attractive for privacy- and governance-sensitive applications, but it also makes aggregation a security-critical operation. At each round, the server must update the global model from client messages whose underlying data and training process cannot be directly inspected. As a result, the aggregation rule becomes the main line of defense when some clients submit harmful updates.

Model-poisoning attacks exploit this trust bottleneck by replacing benign client updates with crafted messages. These attacks can be targeted, such as backdoor and model-replacement attacks \citep{bagdasaryan2020backdoor,xie2020dba,wang2020tails}, or untargeted, where the goal is to degrade clean test performance and destabilize collaborative training \citep{baruch2019alie,fang2020local,shejwalkar2021manipulating}. This paper focuses on the untargeted setting. The core challenge is not only that malicious updates can be arbitrary, but also that benign updates are naturally diverse under non-i.i.d.\ client data. A geometrically unusual update may therefore be either an attack or a legitimate consequence of client heterogeneity \citep{zhao2018noniid,li2020fedprox,karimireddy2020scaffold,zhu2021noniidsurvey}.

Existing Byzantine-robust defenses address this threat through geometric selection, coordinate-wise robust statistics, or auxiliary validation, but their trust signals can be too coarse for heterogeneous FL. Krum and MultiKrum \citep{blanchard2017krum} select centrally located updates, Bulyan \citep{elmhamdi2018bulyan} combines selection with a second robust aggregation stage, coordinate-wise median \citep{yin2018median} suppresses marginal outliers, and RFA \citep{pillutla2022rfa} uses a geometric-median objective. These rules can confuse benign dispersion with malicious deviation when client data are non-i.i.d. Auxiliary-information defenses partly address this limitation, but FLTrust \citep{cao2021fltrust} compresses trusted data into a single reference direction and FLDetector \citep{zhang2022fldetector} relies on historical behavior prediction. Such signals may miss attacks whose harm is concentrated on a subset of classes or semantic regions. The resulting open problem is how to validate client updates by their task-level usefulness rather than only by their location in update space.

We propose \textbf{CARAT}, a class-aware robust aggregation rule that validates updates through hidden, class-balanced tasks sampled from a small labeled server reference set. The key idea is that a benign non-i.i.d.\ update need not be geometrically central, but it should provide useful evidence on semantic validation tasks that are hidden from clients. CARAT first places candidate updates in a common scoring geometry through robust clipping and common-radius evaluation. It then measures task-level loss reduction to form hidden-task certificates, penalizes coordinate-rank anomalies, and solves a capped-simplex weight optimization problem with temporal regularization. The final update is computed from clipped client updates, so the server can preserve useful benign diversity while limiting updates that are semantically harmful or structurally extreme. This design explicitly targets settings where a modest labeled server reference set is available; it is not intended to be a fully data-free defense.

We evaluate CARAT on CIFAR-100~\citep{krizhevsky2009learning} with ResNet18~\citep{he2016deep} under five untargeted model-poisoning attacks, three Dirichlet non-i.i.d.\ regimes, and a $20\%$ malicious-client ratio. In the completed five-seed sweep, CARAT is competitive for $\alpha\in\{1,0.5\}$: under the four conventional attacks it has average rank $2.5$ and remains within $0.1$--$1.2$ percentage points of the strongest row. It also ranks second among defenses with Mimic results. This behavior does not extend uniformly to $\alpha=0.1$, where FangAttack causes four of five seeds to finish at chance accuracy and MinSum causes one similar collapse. The component study further shows only modest separation among variants. We therefore make a bounded claim: task-conditioned validation is a viable reference-assisted signal in moderate heterogeneity, while the current design is unstable under some attacks in the most heterogeneous tested setting. Our main contributions are as follows.

\begin{itemize}
\item We introduce hidden-task certificates for robust FL aggregation, using class-balanced server-side validation tasks to evaluate whether each candidate update improves semantic slices of the learning problem.
\item We develop the CARAT aggregation objective, which combines hidden-task utility, coordinate-rank anomaly penalties, capped per-client influence, and temporal regularization into a single robust weighting rule.
\item We provide an auditable five-seed CIFAR-100/ResNet18 evaluation against public robust-aggregation baselines under five attacks and three non-i.i.d.\ partitions, with explicit reporting of both competitive regimes and observed failure cases.
\end{itemize}
